\section*{Chapter \ref{ch:b3:measure} --- Measure Theory}

\begin{solution}{exo:b3:measure:1}
(a) Complementation swaps the two defining cases. A countable
union of countable sets is countable; if one member is
co-countable, the union is co-countable: stability holds. It
contains the singletons, and any $\sigma$-algebra containing
them contains all countable sets (countable unions) and their
complements: it is $\sigma(\{\text{singletons}\})$.

(b) Write $\mathcal B = \sigma(\text{opens})$. Every open
subset of $\R$ is a countable union of open intervals with
rational endpoints (around each rational point of the open set,
a rational-radius interval inside it), so opens $\in
\sigma(\text{open intervals}) \subseteq \sigma(\text{rational
data})$. Conversions: $\intoo ab = \bigcup_n\intcc{a +
\frac1n}{b - \frac1n}$; $\intcc ab = \bigcap_n \intoo{a -
\frac1n}{b + \frac1n}$; $\intoc{-\infty}a = \bigcap_n
\intoo{-\infty}{a + \frac1n}$ and conversely $\intoo ab =
\intoo{-\infty}b \setminus \intoc{-\infty}a$; rational rays:
$\intoc{-\infty}a =
\bigcap_{q \in \Q,\, q > a}\intoc{-\infty}q$. Each family
generates the others by countable operations: all four generate
$\mathcal B$.

(c) With finite endpoints only, the family is not even an
algebra: the complement of $\intoc01$ contains unbounded rays.
Allowing infinite endpoints ($\intoc{-\infty}b$, $\intoo
a{+\infty}$) it becomes an algebra (complements and finite
unions of such unions are such), but not a $\sigma$-algebra:
$\{0\} = \bigcap_n\intoc{-\frac1n}0$ is not a finite union of
nondegenerate intervals.
\end{solution}

\begin{solution}{exo:b3:measure:2}
(a) $A \cup B = A \sqcup (B \setminus (A\cap B))$, so
$\mu(A\cup B) = \mu(A) + \mu(B) - \mu(A\cap B)$ (finiteness
permits the subtraction). Three sets: apply the two-set formula
twice,
\[
\mu(A\cup B\cup C) = \sum\mu(A) - \sum\mu(A\cap B) +
\mu(A\cap B\cap C)
\]
(sums over the obvious index sets).

(b) For Lebesgue measure, $A_n = [n, +\infty) \downarrow
\varnothing$, but $\lambda(A_n) = \infty \not\to 0$.

(c) A point lies in an interval of length $\varepsilon$:
$\lambda(\{x\}) = 0$; countable subadditivity kills countable
sets. Hence $\lambda(\Q \cap \intcc01) = 0$ and, by additivity,
$\lambda(\intcc01\setminus\Q) = 1$: the irrationals carry all
the length.
\end{solution}

\begin{solution}{exo:b3:measure:3}
The rays $\intoc{-\infty}t$ form a $\pi$-system (the
intersection of two is the smaller) generating $\mathcal B(\R)$
(\cref{exo:b3:measure:1}). The sets $P_k = \intoc{-\infty}k$
increase to $\R$ with $\mu(P_k) \leq 1 < \infty$:
\cref{thm:b3:measure:uniqueness} applies, and $\mu = \nu$ on
$\mathcal B(\R)$. Thus the distribution function determines the
measure.
\end{solution}

\begin{solution}{exo:b3:measure:4}
For every $N$, $\limsup A_n \subseteq \bigcup_{n \geq N}A_n$,
so $\mu(\limsup A_n) \leq \sum_{n\geq N}\mu(A_n)$, the tail of
a convergent series: let $N \to \infty$. Application: with
$A_n = \{x \in \intcc01 : \abs{x - r_n} \leq 4^{-n}\}$
($r_n$ the $n$-th rational), $\lambda(A_n) \leq 2\cdot4^{-n}$
is summable: $\lambda(\limsup A_n) = 0$, i.e.\ almost every $x$
belongs to only finitely many $A_n$. (Yet every $x$ is a limit
of rationals: the point is the \emph{speed} $4^{-n}$.)
\end{solution}

\begin{solution}{exo:b3:measure:5}
$K = \bigcap K_n$ is an intersection of finite unions of closed
intervals: compact. At stage $n$ there remain $2^n$ intervals
of common length $\ell_n \leq 2^{-n}$ (each stage halves and
shrinks); an interval $I \subseteq K$ would lie inside a single
stage-$n$ interval for every $n$, forcing $\ell(I) = 0$: empty
interior. The measure removed is $\sum_{n\geq1}2^{n-1}\cdot
4^{-n} = \frac12\sum_{n\geq1}2^{-n} = \frac12$, all removals
being disjoint open intervals: $\lambda(K) = \frac12$.

Variant: removing central intervals of length $\varepsilon
4^{-n}$ leaves a nowhere dense compact $K^{(\varepsilon)}$ of
measure $1 - \frac\varepsilon2$. Then $\bigcup_m
K^{(1/m)}$ is meagre (countable union of nowhere dense sets)
of measure $\geq \sup_m(1 - \frac1{2m}) = 1$: a meagre set of
full measure --- and its complement in $\intoo01$ is a dense
$G_\delta$ of measure $0$ (topologically fat, metrically
null). The complement of $K^{(\varepsilon)}$ in $\intoo01$ is
open, dense, of measure $\frac\varepsilon2 < \varepsilon$.
\end{solution}

\begin{solution}{exo:b3:measure:6}
Cutting $\intoc01$ into $2^n$ translates of $\intoc0{2^{-n}}$:
$c = 2^n\,\mu(\intoc0{2^{-n}})$, so $\mu(\intoc0{2^{-n}}) =
c\,2^{-n} = c\,\lambda(\intoc0{2^{-n}})$. By translation
invariance and additivity, $\mu = c\lambda$ on every interval
$\intoc ab$ with $b - a$ a dyadic rational and any $a$; a
general $\intoc ab$ is an increasing union of such
($b_k \uparrow b$ dyadic steps from $a$), and continuity from
below extends the equality. The intervals $\intoc ab$ form a
$\pi$-system generating $\mathcal B(\R)$, with $\intoc{-k}k
\uparrow \R$ of finite measure ($\mu(\intoc{-k}k) = 2kc$):
\cref{thm:b3:measure:uniqueness} gives $\mu = c\lambda$ on
$\mathcal B(\R)$.
\end{solution}

\begin{solution}{exo:b3:measure:7}
By regularity (\cref{thm:b3:measure:regularity}) choose $K
\subseteq A \subseteq U$ with $K$ compact, $U$ open,
$\lambda(U\setminus K) < \varepsilon$ (both approximations
within $\varepsilon/2$, and $\lambda(U \setminus K) =
\lambda(U\setminus A) + \lambda(A \setminus K)$). Write $U$ as
a countable disjoint union of open intervals $(I_n)$ (the
components of the open set); the compact $K$ is covered by
finitely many, $K \subseteq B = I_1\cup\dots\cup I_N \subseteq
U$. Then $A \setminus B \subseteq A\setminus K \subseteq
U\setminus K$ and $B \setminus A \subseteq U \setminus A
\subseteq U\setminus K$:
$\lambda(A\,\triangle\,B) \leq 2\lambda(U\setminus K)$ ---
start from $\varepsilon/2$ to land below $\varepsilon$.
\end{solution}

\begin{solution}{exo:b3:measure:8}
Replacing $A$ by $A \cap [-M, M]$ of positive measure (some $M$
works, by continuity from below), assume $0 < \lambda(A) <
\infty$. Take $U \supseteq A$ open with $\lambda(U) <
\frac43\lambda(A)$ and decompose $U = \bigsqcup_nI_n$ into
disjoint open intervals: $\lambda(A) = \sum_n\lambda(A\cap
I_n)$. If every $n$ had $\lambda(A\cap I_n) \leq
\frac34\ell(I_n)$, summing would give $\lambda(A) \leq
\frac34\lambda(U) < \lambda(A)$: some interval $I$ satisfies
$\lambda(A\cap I) > \frac34\ell(I)$. Set $B = A \cap I$ and let
$\abs t < \frac12\ell(I)$: both $B$ and $B + t$ lie in the
interval $I \cup (I + t)$, of length $< \frac32\ell(I)$. If
they were disjoint: $\lambda(B) + \lambda(B + t) = 2\lambda(B)
> \frac32\ell(I)$ would exceed the containing interval's
measure --- impossible. So $B \cap (B + t) \neq \varnothing$:
some $x \in B$ writes $x = y + t$ with $y \in B$, and $t = x -
y \in A - A$.
Hence $\intoo{-\frac{\ell(I)}2}{\frac{\ell(I)}2} \subseteq A -
A$.
\end{solution}

\begin{solution}{exo:b3:measure:9}
The Vitali translates $(V + q)_{q\in\Q}$ partition $\R$ (every
real is equivalent to exactly one representative). Suppose all
the sets $B_q = A \cap (V + q)$ were measurable. Any two
elements of $V + q$ differ by an irrational or zero (two
distinct representatives are inequivalent), so $B_q - B_q$
meets $\Q$ only in $\{0\}$: it contains no interval, and
Steinhaus (\cref{exo:b3:measure:8}) forces $\lambda(B_q) = 0$.
Then $\lambda(A) \leq \sum_{q}\lambda(B_q) = 0$, contradicting
$\lambda(A) > 0$. So some $B_q \subseteq A$ is non-measurable.
\end{solution}

\begin{solution}{exo:b3:measure:10}
\emph{Measurable $\Rightarrow$ $\varepsilon$-approximation}: by
outer regularity (\cref{thm:b3:measure:regularity}) pick open
$U \supseteq A$ with $\lambda(U) \leq \lambda(A) +
\varepsilon$; measurability allows the subtraction
$\lambda(U\setminus A) = \lambda(U) - \lambda(A) \leq
\varepsilon$. \emph{$\varepsilon$-version $\Rightarrow$
$G_\delta$-version}: take $U_n$ with $\lambda^*(U_n\setminus
A) < \frac1n$ and $G = \bigcap U_n$: a $G_\delta$ with
$\lambda^*(G\setminus A) \leq \lambda^*(U_n\setminus A) \to 0$.
\emph{$G_\delta$-version $\Rightarrow$ measurable}: $G\setminus
A$ is $\lambda^*$-null, hence measurable by completeness
(\cref{thm:b3:measure:caratheodory}), and $A = G \setminus
(G\setminus A)$ is measurable ($G$ is Borel). So Lebesgue sets
are exactly ``Borel modulo null''.
\end{solution}

\begin{solution}{pb:b3:measure:1}
\textbf{1.} Induction. Continuity: the three formulas agree at
the junctions ($\frac12c_n(1) = \frac12$ and $\frac12 +
\frac12c_n(0) = \frac12$); each piece is continuous.
Monotonicity and the boundary values are inherited. For the
contraction estimate: on $[0,\frac13]$, $\abs{c_{n+1} - c_n}(x)
= \frac12\abs{c_n - c_{n-1}}(3x) \leq \frac12\norm{c_n -
c_{n-1}}_\infty$; on the middle third the difference is $0$; on
the right third, the same as the left.

\textbf{2.} $\norm{c_{n+1} - c_n}_\infty \leq
2^{-n}\norm{c_1 - c_0}_\infty$: the series of increments
converges uniformly, so $c_n \to c$ uniformly; $c$ is
continuous, nondecreasing, $c(0) = 0$, $c(1) = 1$ (all
preserved by uniform limits), and passing to the limit in the
defining recursion shows $c$ itself satisfies the three
self-similar identities.

\textbf{3.} By the middle identity, $c \equiv \frac12$ on
$\intcc{\frac13}{\frac23}$, the first gap. Every gap of $C$ is
the image of the first gap under a composition of the two
affine contractions $x \mapsto \frac x3$, $x\mapsto\frac{x +
2}3$; the identities transport constancy accordingly (with
values the dyadic rationals). For the digit formula, take $x =
\sum_n 2b_n3^{-n} \in C$: if $b_1 = 0$ then $x \in
[0,\frac13]$ and $c(x) = \frac12c(3x)$ with $3x$ having digits
$(b_2, b_3, \dots)$; if $b_1 = 1$ then $x \in [\frac23, 1]$
and $c(x) = \frac12 + \frac12c(3x - 2)$, same shift. By
induction, the first $N$ binary digits of $c(x)$ are $b_1,
\dots, b_N$ for every $N$: $c(x) = \sum_nb_n2^{-n}$.

\textbf{4.} Off $C$, $c$ is locally constant: differentiable
with derivative $0$. Since $\lambda(C) = 0$
(\cref{ex:b3:measure:cantor}), $c' = 0$ almost everywhere. Yet
$c(1) - c(0) = 1$: the fundamental theorem in its $\mathcal
C^1$ form requires $c$ to be differentiable \emph{everywhere}
with continuous (or at least integrable, plus absolute
continuity --- see \cref{ch:b3:lebesgue}) derivative; $c$ is
not differentiable at points of $C$, and more fundamentally $c$
fails \emph{absolute continuity}: it climbs on a null set.

\textbf{5.} Given $y = \sum_n\beta_n2^{-n} \in \intcc01$
($\beta_n \in \{0,1\}$), the point $x = \sum_n2\beta_n3^{-n}
\in C$ has $c(x) = y$ by question 3: $c(C) = \intcc01$, a set
of measure $1$ --- the null set $C$ carries, through $c$, the
whole interval.

\textbf{6.} $h$ is continuous, and strictly increasing ($x$ is,
$c$ is nondecreasing); $h(0) = 0$, $h(1) = 1$, so by the
intermediate value theorem $h$ is a continuous bijection of
$\intcc01$; a continuous bijection from a compact space to a
Hausdorff one is a homeomorphism
(\cref{cor:b3:topology:compacthomeo}).

\textbf{7.} On a gap $(u, v)$ (length $\ell$), $c$ is constant,
so $h$ is affine with slope $\frac12$: $h((u,v))$ is an
interval of length $\frac\ell2$. The gaps are disjoint and $h$
is injective: the images are disjoint, of total measure
$\frac12\sum\ell = \frac12(1 - \lambda(C)) = \frac12$.

\textbf{8.} $h(\intcc01) = \intcc01$ and $h(C)$ is compact
(continuous image), hence measurable, with
\[
\lambda\bigl(h(C)\bigr) = 1 -
\lambda\bigl(h(\intcc01\setminus C)\bigr) = 1 - \tfrac12 =
\tfrac12 .
\]
A homeomorphism can inflate a null set to measure $\frac12$:
``topological size'' (category, dimension) and ``measure'' are
transported by homeomorphisms very differently --- only the
former is a topological invariant.

\textbf{9.} $Z = h^{-1}(W) \subseteq h^{-1}(h(C)) = C$ has
$\lambda^*(Z) \leq \lambda(C) = 0$: a null set, hence
Lebesgue-measurable by completeness
(\cref{thm:b3:measure:caratheodory}).

\textbf{10.} Let $\varphi$ be continuous and $\mathcal D = \{B
: \varphi^{-1}(B) \in \mathcal B\}$. Preimages commute with
complements and countable unions, so $\mathcal D$ is a
$\sigma$-algebra; it contains the open sets (continuity):
$\mathcal D \supseteq \mathcal B$ --- preimages of Borel sets
under continuous maps are Borel.

\textbf{11.} If $Z$ were Borel, apply question 10 to the
continuous $\varphi = h^{-1}$: $\varphi^{-1}(Z) = h(Z) = W$
would be Borel, hence Lebesgue-measurable --- contradicting the
choice of $W$. So $Z \in \mathcal L \setminus \mathcal B$:
Lebesgue's $\sigma$-algebra strictly contains Borel's.

\textbf{12.} $g = \mathbf 1_Z$ is Lebesgue-measurable ($Z \in
\mathcal L$) and $\varphi = h^{-1}$ is continuous, but
$(g\circ\varphi)^{-1}(\{1\}) = \varphi^{-1}(Z) = W$ is not
measurable: $g \circ \varphi$ is not Lebesgue-measurable. The
care needed: ``Lebesgue-measurable function'' means preimages
of \emph{Borel} sets land in $\mathcal L$; composing requires
preimages of \emph{Lebesgue} sets to be Lebesgue, which
continuity does not grant (here $\varphi^{-1}(Z) \notin
\mathcal L$ even though $\varphi$ is a homeomorphism).

\textbf{13.} The chains, with witnesses for strictness:
\[
\{\text{countable}\} \subsetneq \{\text{Borel null}\}
\subsetneq \{\text{Lebesgue null}\} \subsetneq \mathcal L
\subsetneq \mathcal P(\R),
\qquad
\{\text{Borel null}\} \subsetneq \mathcal B \subsetneq
\mathcal L .
\]
Witnesses: $C$ is Borel, null, uncountable (first gap); $Z$ is
Lebesgue null but not Borel (second and, inside $\mathcal B
\subsetneq \mathcal L$, sixth); a fat Cantor set is Borel,
nowhere dense, of positive measure (separating null sets from
Borel sets); Vitali's $V$ is not in $\mathcal L$ (last gap).
Measure, topology and cardinality slice $\mathcal P(\R)$ along
genuinely different lines.
\end{solution}
